\chapter{Conclusions and Future Work}
\label{chap:conclusions-future-work}

This chapter concludes the thesis by consolidating its principal contributions, practical implications, limitations, and future extensions. The research addressed the central question of how predictive operational analytics can improve arrival management at Blaise Diagne International Airport. It responded to this question through an arrival-centric predictive operational intelligence framework that connects operational analysis, short-horizon forecasting, discrete-event simulation, key performance indicators, and manager-facing decision support.

The study was motivated by the operational pressure created when multiple flights arrive in short succession, generating concentrated demand across immigration, baggage reclaim, customs, and terminal exit processes. In this environment, operational monitoring alone is insufficient because it primarily describes existing conditions. Effective arrival management also requires the ability to anticipate passenger and baggage pressure, evaluate the likely consequences of that pressure, and communicate the resulting information in a form that supports timely operational action.

The thesis, therefore, followed the analytical sequence established in the Introduction: explain, predict, and support decisions. Operational and temporal patterns were first analyzed to explain the development of passenger pressure, baggage pressure, and congestion risk. Predictive models were then used to estimate short-horizon arrival-side pressure. Finally, the forecasts were integrated with AnyLogic simulation metrics and Power BI dashboards to support resource planning and operational decision-making. The resulting framework demonstrates how arrival data can be transformed from a descriptive record of airport activity into operational intelligence that supports a more proactive approach to congestion management.

\section{Summary of Contributions}
\label{sec:summary-contributions}

The first contribution of the thesis is the development of an arrival-centric predictive operational intelligence framework. The framework was designed around the operational characteristics of the AIBD arrival process rather than around airport activity in general. Its analytical scope therefore reflects the sequence through which arriving passengers move from flight arrival to immigration, baggage reclaim, customs, and terminal exit.

This focus is necessary because pressure within the arrival system does not arise from a single isolated activity. The simulation results reported in \Cref{sec:simulation-results} demonstrate the interdependence empirically: increasing immigration-processing capacity in Scenario A reduced the immigration queue from 80 to 5 passengers, a 94 percent reduction, but simultaneously increased the baggage queue from 120 to 198 passengers and produced no improvement in average passenger system time, which rose marginally from 102 to 106 minutes. Congestion was relocated rather than removed. The framework addresses this interdependence by connecting passenger and baggage indicators with forecasting, simulation, performance measurement, and decision support, so that the consequences of a single-subsystem intervention can be assessed across the full arrival chain before implementation.

The second contribution is the establishment of validated passenger- and baggage-pressure indicators. The research was grounded in a validated master dataset of 22,270 commercial arrival records, derived from an initial FR24 archive of 50,385 flight movements through the documented five-step validation and filtering pipeline reported in \Cref{tab:data-validation-pipeline}. The indicators provided a consistent analytical basis across the components of the thesis: the same pressure definitions were reproduced independently as Power BI DAX-calculated columns, ensuring that the operational categories displayed to managers are identical to those used in the Python analytics and in the simulation input tables.

Passenger pressure and baggage pressure were treated as related but distinct dimensions of arrival operations. Passenger pressure represents the demand generated by arriving passengers across immigration, baggage reclaim, customs, and exit processes. It is classified into four operational bands using fixed thresholds of 230, 190, and 160 average passengers per hour. Baggage pressure represents the corresponding demand placed on baggage-handling and reclaim capacity, classified into three bands using the 33rd and 66th percentiles of the hourly estimated-bags series. The distinction is operationally relevant because baggage reclaim imposes a capacity constraint distinct from immigration processing, and, on this basis, the two are modeled separately in the simulation. It is not, however, a distinction between two independent demand signals: as established in \Cref{sec:analytics}, \nolinkurl{estimated_bags} is derived from \nolinkurl{estimated_passengers} by a fixed multiplier, so the two series are identical in shape. The wider baggage ``High'' band (02:00, 04:00--06:00, 16:00--17:00, and 20:00--21:00) follows from the percentile-based thresholds applied to baggage rather than from any difference in demand behavior. It must not be read as evidence that baggage pressure persists beyond the passenger surges that generate it.

The empirical results obtained from these indicators indicate two recurring daily arrival banks at AIBD. The morning bank spans approximately 04:00--06:00 and peaks at 203.4 estimated passengers at 05:00; the evening bank spans approximately 19:00--21:00 and peaks at 210.8 estimated passengers at 20:00, the single highest-pressure hour of the operational day. Between these banks, every hour from 07:00 to 15:00 is classified as ``Low'' pressure, with a minimum of 135.7 passengers at 11:00, 36 percent below the evening peak. These two windows constitute the empirical basis for the 160-passenger operational threshold applied throughout the framework. They are the specific periods to which the practical implications in \Cref{sec:practical-implications} refer.

The third contribution is the integration of forecasting, simulation, and decision support within a single operational framework, in which each component performs a specific function within a connected decision process rather than standing as an independent analytical exercise. Forecasting provided short-horizon estimates of passenger pressure and congestion conditions, an emphasis consistent with the finding of \textcite{orsini2019neural} that terminal passenger-flow forecasts are most reliable at horizons of approximately twenty minutes or less. Four regression algorithms were evaluated for passenger-pressure prediction, with LightGBM achieving the best performance (MAE 34.06, RMSE 46.61, $R^2$ 0.2366) and XGBoost retained as the saved production model; four classifiers were evaluated for congestion detection, with the reduced-feature LightGBM model selected for deployment based on its accuracy (0.7805), F1 score (0.4690), and ROC-AUC (0.7914). Their value was therefore assessed not only through predictive performance but also through operational interpretability.

A methodological contribution arises within this component. The initial full-feature classifiers returned near-perfect scores, with three of the four models reporting accuracy, precision, recall, F1, and ROC-AUC of 1.0000. Rather than reporting these values as evidence of predictive capability, the study identified the retained \nolinkurl{future_baggage_pressure} variable as a source of information leakage; its Random Forest importance of approximately 0.60 exceeded that of all other predictors combined, and the classifiers were retrained on a reduced 384-feature specification. The resulting ROC-AUC values of 0.7677--0.7973 reflect the model’s genuine operational discrimination. The identification and correction of this leakage, and the decision to report both specifications transparently, are themselves contributions to the methodological rigor of the framework.

The AnyLogic discrete-event simulation extended the analysis from predicted demand to potential operational consequences by representing passenger generation and movement through immigration, baggage reclaim, customs, and exit, as well as queue accumulation and congestion propagation, over a 12-hour operational horizon spanning 720 simulation minutes. Six named scenarios covering the baseline, capacity-optimization, peak-stress, dynamic-intervention, and integrated-optimization configurations were compared through the specified key performance indicators: passenger system time, the three process-specific queue sizes, throughput, congestion severity, and operational resilience. The consolidated KPI outcomes are reported in \Cref{tab:simulation-kpi-results}.

The integration of predictive indicators with these metrics connected the identification of future pressure to an assessment of its possible operational effects. Two findings emerge with particular clarity. Scenario B (Baggage Optimization) produced the highest throughput of the six scenarios (573 exited passengers), the lowest total queue (145 passengers), and the highest \nolinkurl{Resilience_Score} (0.76), indicating that baggage-reclaim capacity is the dominant operational constraint under baseline conditions. Scenario E (Integrated Optimization) produced the lowest average system time (78 minutes, a 24 percent reduction relative to the 102-minute baseline) and the second-highest \nolinkurl{Resilience_Score} (0.75), indicating that coordinated multi-subsystem improvement outperforms isolated intervention on the passenger-experience measure, on the single run recorded for each Scenario.

The fourth contribution is the development of manager-facing predictive intelligence. The Power BI component completed the analytical sequence by translating forecasts, simulation results, pressure classifications, and key performance indicators into decision-support outputs, organized into operational overview, scenario comparison, dynamic intervention monitoring, executive decision support, and operational intelligence insights. The research thereby addressed the gap between producing predictions and supporting operational decisions: a forecast has limited practical value when its meaning cannot be interpreted in relation to airport capacity, congestion risk, resource requirements, or possible interventions.

Within these pages, the 160-passenger operational threshold is displayed as an explicit reference line against the hourly passenger-pressure forecast, and the scenario-comparison ranking matrix assigns each simulated configuration a qualitative risk level and recommendation status (\Cref{sec:decision-support-system}).

The Key Influencers results provide independent corroboration of the framework’s central finding. The analysis identified Passengers Pressure Level being High as the strongest single influencer of high congestion risk, associated with a 3.17-fold increase in likelihood, followed by the arrival hour falling between 19:00 and 21:00 (3.02-fold), Baggage Pressure Level being High (2.23-fold), and the arrival hour being 06:00 or earlier (1.48-fold). Because this analysis applies Power BI’s own internal statistical procedure, it was computed independently of the machine-learning classifiers and the SHAP attribution reported in Chapter 5. The convergence of all three procedures on the same two temporal drivers, the evening and morning arrival banks, constitutes methodologically independent confirmation that time-of-day is the dominant determinant of congestion risk at AIBD.

The fifth contribution is the framework’s AIBD-specific operational relevance. AIBD’s passenger growth, strategic regional position, distance from Dakar, and multi-agency arrival environment create demands that cannot be addressed solely through infrastructure; they also require timely operational information and coordination among the organizations responsible for processing arriving passengers. The framework responds by focusing on periods when several flights converge and create concentrated pressure, and by treating immigration, baggage reclaim, and customs as operationally connected, even though they involve different stakeholders and resources. This orientation aligns directly with the strategic priorities articulated by AIBD leadership, who identified AI-powered systems, predictive analytics, crowd management optimization, smart queuing systems, and dynamic staff allocation as means of forecasting peak passenger flows and optimizing staffing allocation \parencite{demir2025speaker}.

The AIBD-specific contribution also concerns the relationship between technological modernization and passenger experience. The thesis does not position predictive intelligence as an alternative to human-centered service. It positions data-driven decision-making as a means of supporting the efficient, professional, and welcoming arrival experience associated with Senegalese Teranga. Reducing uncertainty, anticipating peaks, and preparing resources can support both operational performance and the quality of interaction experienced by arriving passengers.

The framework's contribution can also be positioned in relation to the five gaps identified in \Cref{chap:literature-review}. It addresses the absence of arrival-centric studies by modelling the connected immigration-baggage-customs-exit sequence rather than a single process; it addresses the limited attention to the baggage side of the arrival chain by modelling baggage-reclaim capacity as a distinct operational constraint with its own classification thresholds, while establishing that the baggage series derived here is not an independent demand signal; it addresses the absence of integrated frameworks by linking forecasting, simulation, KPI evaluation and dashboard reporting within one continuous pipeline; it addresses the absence of manager-facing predictive intelligence by translating model outputs into pressure classifications, scenario rankings and operational recommendations; and it addresses the absence of African airport research by developing and evaluating the framework at a West African airport, a context absent from all eight studies reviewed in \Cref{chap:literature-review}.

Taken together, the contributions establish a coherent process for explaining, predicting, evaluating, and communicating arrival pressure. The framework combines validated indicators, short-horizon forecasting, scenario-based simulation, operational key performance indicators, and manager-facing dashboards. In doing so, it converts arrival management from a predominantly reactive activity into a more proactive, data-driven decision-making process.

\section{Practical Implications}
\label{sec:practical-implications}

The practical implications of the research concern the use of predictive operational intelligence to prepare airport resources before pressure develops into severe congestion. The framework does not replace the responsibilities of airport managers, immigration officers, airline operations teams, baggage-handling personnel, or other stakeholders. Rather, it provides these stakeholders with structured information that can support the timing and coordination of their decisions.

The first implication concerns staffing preparation. Arrival pressure is most difficult to manage when several flights converge within a limited time window, and staffing is not aligned with the resulting demand. Without short-horizon forecasts, the need for additional capacity may be identified only after queues have accumulated. Passenger-pressure forecasts provide advance information on the timing and relative intensity of that demand, and this study identifies the specific windows to which preparation should be directed: approximately 04:00--06:00 and 19:00--21:00, with 20:00 representing the highest single-hour demand of the operational day at 210.8 estimated passengers. The practical contribution lies in matching preparation to these two windows rather than distributing staffing uniformly across an operational day, during which nine consecutive hours (07:00--15:00) consistently fall below the 160-passenger operational threshold.

For immigration officers, reliable information about expected peaks can support counter planning and queue management. For airport managers, the same information can support coordination across the organizations involved in the arrival process. The Scenario A result reported in \Cref{sec:simulation-results} provides a specific caution in this respect: increasing immigration staffing alone reduced the immigration queue by 94 percent but produced no reduction in average passenger system time, because the relieved demand accumulated instead at baggage reclaim. The framework, therefore, treats staffing as a coordinated operational requirement rather than as an isolated decision within a single processing stage, and provides the quantitative basis for that coordination.

The second implication concerns baggage-capacity planning. Passenger and baggage demand are closely connected, but baggage handling and reclaim introduce their own capacity requirements, so pressure may develop at reclaim facilities even where passenger processing at other stages remains manageable. The baggage-pressure indicators identify the periods in which reclaim demand is highest, reaching a maximum of 421.5 estimated bags at 20:00. Because the baggage series is a fixed multiple of the passenger series (\Cref{sec:analytics}), those periods coincide with the two principal arrival banks; the additional hours classified ``High'' under the percentile-based baggage thresholds (02:00 and 16:00--17:00) follow from the choice of classification rule rather than from independent baggage behaviour. They should be read as a property of the indicator rather than of the demand. Integration with the simulation metrics enables assessment of the resulting pressure against baggage queue size, passenger system time, throughput, and congestion severity.

The practical importance of this capability extends beyond the baggage-reclaim process itself. Baggage delays can increase the time passengers spend in the arrival system and may lead to passenger accumulation before customs and at the terminal exit. Baggage-capacity planning is therefore connected to the performance of the wider arrival chain. The Scenario B result quantifies this connection: raising the baggage-capacity multiplier from 1.00 to 1.30 reduced the baggage queue from 120 to 46 passengers, increased throughput from 484 to 573 exited passengers, and reduced average system time from 102 to 83 minutes, a system-wide improvement of 19 minutes obtained through a single-subsystem intervention. Among the interventions tested, this represents the highest operational return and identifies baggage-reclaim capacity as the priority investment target.

The third implication concerns queue-management strategies. Immigration, baggage reclaim, and customs can each become bottlenecks when demand exceeds available processing capacity. Because these processes are sequentially connected, queue growth at one stage may affect those that follow. The simulation framework evaluates queue accumulation under the baseline, capacity-optimization, peak-stress, dynamic-intervention, and integrated-optimization scenarios. Its results demonstrate that the relevant indicators do not move together: Scenario A improved one queue measure dramatically while worsening another, and Scenarios B and E each traded a reduction in one queue for a partial increase in the other. Evaluating a queue-management decision against a single stage-specific measure would therefore produce a misleading assessment of its system-level effect. The framework makes these consequences visible before they are encountered in actual operations, so that queue management can be approached as a system-level issue rather than as a set of separate responses to immigration, baggage, or customs congestion.

The fourth implication concerns intervention timing. Whether an operating model is reactive or proactive depends substantially on when information becomes available and when a response is initiated; an intervention introduced after queues have reached a critical level may be less effective than preparation undertaken before the peak in arrivals. Short-horizon forecasting supports this temporal requirement by identifying anticipated pressure before its operational effects have fully developed. The Scenario D result reported in \Cref{sec:simulation-results} establishes an important boundary on this capability: the queue-triggered intervention rules activated during execution produce KPI outcomes numerically identical to those of the equivalent non-adaptive Scenario C configuration under a +50 percent arrival surge. This indicates the existence of an operational saturation limit beyond which reactive capacity adjustment, however well timed, cannot compensate for a demand-capacity shortfall, and reinforces the case for preparation in advance of the peak rather than adjustment during it.

The fifth implication concerns the multi-stakeholder benefits of improved arrival management. Arrival pressure is not solely a technical problem related to queues or processing capacity; it affects passengers, airport staff, immigration and customs personnel, airlines, ground handlers, and airport management. Citizens, tourists and business travellers benefit from reduced uncertainty and more predictable processing and baggage-reclaim times; immigration officers benefit from advance information concerning passenger peaks, allowing staffing and queue arrangements to be prepared in relation to expected demand; airline operations and ground-handling teams benefit from improved coordination of baggage processing; and airport managers benefit from an integrated view of pressure, capacity, simulation outcomes and recommended operational attention. These benefits are interdependent: improvements in staffing preparation reduce queue accumulation; better baggage-capacity planning reduces passenger system time and downstream congestion; and more consistent arrival performance supports airline service quality and reinforces AIBD’s role as Senegal’s principal international gateway.

The framework thus supports both operational efficiency and human-centered service. This relationship is particularly relevant to AIBD’s objective of combining technological modernization with Teranga, as well as to the airport leadership’s stated intention to deploy e-gates, self-service technologies, and intelligent queue management as part of operational modernization \parencite{demir2025speaker}. The immigration-capacity scenarios evaluated in this thesis provide a quantitative basis for anticipating the downstream baggage-side consequences of such immigration-side investments before implementation. Predictive operational intelligence contributes to this objective by supporting preparation, reducing avoidable congestion, and enabling staff to respond to passenger needs under more controlled operating conditions.

\section{Limitations}
\label{sec:conclusion-limitations}

The findings of the thesis must be understood within the boundaries of the data, analytical models, forecasting purpose, simulation framework, and Decision Support System used in the research. These limitations do not invalidate the framework or its results. They define the scope within which the findings should be interpreted and applied.

The first category concerns data limitations, and a specific consequence follows from the estimation procedure. Because the flight-tracking archive records aircraft movements rather than passenger counts, passenger and baggage demand were derived from a fixed aircraft-type-to-seat-capacity lookup, a constant load factor of 0.82, and the assumption of 2 bags per passenger at an average of 23 kilograms each. These constants were applied uniformly across all routes, airlines, and seasons and were not calibrated against observed passenger manifests or baggage counts. The pressure magnitudes reported throughout the thesis are therefore internally consistent across the analytics, forecasting, simulation and dashboard layers, and are valid for identifying the relative timing and comparative intensity of arrival banks, but they should be interpreted as a consistent proxy for demand rather than as validated absolute measurements.

The second category concerns modeling limitations. The framework depends on the analytical indicators, features, classifications, and models implemented in the study, and these representations were developed to support operational interpretation and short-horizon forecasting. They provide a structured analytical description of arrival pressure but remain dependent on the definitions and modeling logic used in the thesis. The passenger-pressure thresholds of 230, 190, and 160 passengers per hour, and the 180-passenger congestion-classification threshold derived as the 75th percentile of the engineered dataset, are calibrated to AIBD’s observed demand distribution and are not externally derived service-level standards. The results, therefore, reflect the selected operational representation of demand and congestion.

The third category concerns forecasting limitations. The forecasting component was designed for a short-horizon operational purpose, and three specific constraints qualify its results. First, the regression models explain a limited proportion of variance in future passenger pressure, with $R^2$ ranging from 0.1803 to 0.2366 across the four algorithms tested; a substantial component of short-horizon demand variation is therefore not captured by the temporal, cyclical, and operational-history features engineered in this study. Second, the deployed congestion classifiers achieve only moderate recall, ranging from 0.1926 to 0.3662, indicating that a considerable proportion of genuine high-congestion periods would not be flagged in advance under the evaluation conditions tested, an important operational qualification given that early detection was the stated selection priority. Third, the time-based validation returned a perfect 1.0000 across all five metrics despite excluding the direct leakage variables, a result inconsistent with the more moderate random-split reduced-feature performance and one that could not be resolved from the notebooks and exports available. It is reported as an open methodological question rather than as evidence of superior chronological performance.

The classification of congestion periods also depends on the pressure indicators and congestion thresholds implemented in the study. These classifications provide an interpretable basis for identifying operational risk, but their meaning remains connected to the framework in which they were developed. Forecasting accuracy and operational usefulness must therefore be considered together.

The fourth category concerns simulation assumptions, and three limitations specifically apply to the simulation evidence. First, the reported results derive from a single 720-minute run per scenario; the available documentation does not report multiple replications, confidence intervals, or run-to-run variance, so the KPI values should be interpreted as single deterministic outcomes rather than as statistically characterized estimates. This directly affects the interpretation of the identical Scenario C and Scenario D results, for which it cannot be established from the available material whether the equality reflects a genuine operational ceiling or an artifact of the specific runs recorded. Second, the congestion-severity flag reported in the simulation-narrative documentation is uniformly ``CRITICAL'' across all six scenarios, including the best-performing Scenario B and Scenario E, and does not correspond to the graded \nolinkurl{Risk_Level} classification used in the scenario-results table and the Power BI dashboard; the two classifications are computed independently and are not currently reconciled. Third, the four normalized performance scores are computed by max-normalization relative only to the six scenarios tested, and therefore express relative rather than absolute performance; a \nolinkurl{Resilience_Score} of 0.76 indicates the best configuration within this comparison set, not an externally validated benchmark.

Passenger system time, immigration queue size, baggage queue size, customs queue size, throughput, congestion severity, and operational resilience provide a consistent basis for scenario comparison. However, these outputs represent the simulated operational environment. Their primary value lies in supporting structured evaluation of bottlenecks, capacity changes, peak conditions, and interventions within the scope of the implemented model.

The fifth category concerns Decision Support System constraints. The dashboards provide an operational overview, scenario comparison, dynamic-intervention monitoring, executive decision support, and operational intelligence insights. These functions improve the accessibility and interpretability of the results, but the Decision Support System remains constrained by the data, indicators, forecasts, scenarios, and metrics it integrates. The system operates on exported scenario-result tables and historical aggregations rather than on a live operational data feed; it therefore constitutes a decision-support environment demonstrated on historical and simulated data rather than a deployed real-time operational system.

A final limitation concerns the scope of validation. The framework was developed and evaluated as a single-airport case study at AIBD. Its pressure thresholds, congestion classifications, scenario configurations, and resilience values reflect AIBD’s traffic composition, terminal layout, and operational assumptions. The framework has not been evaluated in live operation, and its operational benefits have been demonstrated through simulation rather than through measured improvements in actual airport performance.

Collectively, these limitations establish the conditions under which the thesis findings should be interpreted. The research provides an integrated and operationally relevant framework, but its outputs remain linked to the validated dataset, the defined pressure indicators, the selected forecasting approach, the implemented AnyLogic model, the specified simulation scenarios, and the Power BI structure. These boundaries also provide the basis for the future extensions identified in the following section.

\section{Future Extensions}
\label{sec:future-extensions}

Future work should extend the framework without changing its central operational logic. The relevant directions concern improvements to the pressure indicators, enhancements to the forecasting models, expansion of the simulation scenarios, refinement of the Decision Support System, broader operational integration, and assessment of applicability to other airports.

The first direction is to continue improving the passenger- and baggage-pressure indicators, refining them while preserving their operational purpose. The most direct improvement would be to replace the fixed load-factor and bags-per-passenger constants with observed passenger and baggage counts from airline or ground-handling records, thereby converting the pressure indicators from a consistent proxy into a validated measurement and permitting the absolute magnitude of the reported pressure values to be verified. Such refinement should maintain the connection between flight-level arrival data and the resource demands experienced across immigration, baggage reclaim, customs, and exit, and should remain aligned with operational interpretability so that changes in the indicators continue to have a clear meaning for staffing, baggage capacity, queue management, and intervention timing.

The second direction concerns enhancements to the forecasting models, and three specific extensions follow from the limitations identified in \Cref{sec:conclusion-limitations}. First, the modest regression performance obtained here suggests that sequence-based architectures may capture short-horizon temporal structure more effectively than tree-based ensembles; the direct LSTM strategy identified by \textcite{orsini2019neural} as performing best at short horizons in terminal passenger-flow forecasting provides a specific benchmark against which the present results could be benchmarked. Second, the low recall of the deployed classifiers indicates that class-imbalance treatment, through resampling, class weighting, or decision-threshold adjustment away from the default 0.5, represents a direct route to improving the operational usefulness of congestion detection, since early identification of high-congestion periods was the stated selection priority. Third, the anomalous time-based validation result requires diagnostic re-examination before chronological validation can be relied upon as evidence of deployment-ready performance.

The third direction is the expansion of the simulation scenarios, and the most immediate methodological priority within it is the execution of multiple replications per scenario with reported confidence intervals, which would establish whether the observed differences between scenarios, and in particular the identical Scenario C and Scenario D outcomes, are statistically robust or attributable to single-run variation. Beyond this, the fixed-multiplier scenario design used here could be extended towards the joint prediction-optimization formulations demonstrated in the reviewed literature: the chance-constrained optimization applied by \textcite{cao2025robust} to security-queue management and the machine-learning-enhanced landing-scheduling model of \textcite{pang2024machine} both achieved larger operational gains than rule-based capacity adjustment, and represent a plausible route to improving on the 24 percent system-time reduction obtained by Scenario E. The finding that Scenario D’s queue-triggered rules did not outperform the equivalent static configuration under surge conditions strengthens the case for such a formulation. Expanded scenarios should retain the existing KPI structure to ensure comparability across passenger system time, queue sizes, throughput, congestion severity, and operational resilience.

The fourth direction concerns refining the Power BI Decision Support System. The principal extension would be to connect the dashboard to a live or near-live operational data feed, converting the current demonstration on historical and simulated data into an operational monitoring environment capable of supporting real-time decision-making. A secondary priority is reconciling the simulation model’s internal congestion-severity flag with the graded \nolinkurl{Risk_Level} classification used in the dashboard, ensuring a single, consistent congestion definition applies across the simulation and decision-support layers. Further development should preserve the distinction between descriptive monitoring and predictive decision support, since the value of the system lies in showing not only what has occurred but what pressure is anticipated, how it may affect operational performance, and where managerial attention may be required.

The fifth direction is broader operational integration. Extending the framework beyond the arrival stages examined here should strengthen coordination among the airport operator, immigration authorities, customs services, airlines, ground-handling teams, baggage personnel, and other stakeholders involved in arrival management, with the objective of using shared predictive information to improve preparation and response across interconnected operational responsibilities. The bottleneck-migration results reported in \Cref{sec:simulation-results} provide the operational justification for this direction: because a capacity decision taken by one agency measurably changes the demand experienced by another, shared situational awareness is a functional requirement of effective arrival management rather than an organizational preference.

The sixth direction concerns the applicability of the framework to other airports. The thesis was designed for AIBD and reflects its arrival processes, operational priorities, stakeholder environment, and modernization objectives, so its immediate contribution is AIBD-specific. Future work can assess the extent to which the approach transfers to other settings while preserving its central components: validated arrival data, passenger- and baggage-pressure indicators, short-horizon forecasting, discrete-event simulation, operational key performance indicators, and manager-facing decision support. Given that the reviewed literature contains no studies situated in sub-Saharan or West African airport contexts, extending to comparable regional airports would address the broader geographical gap identified in \Cref{sec:literature-gaps} more effectively than a single case study can.

Application elsewhere would require attention to the airport's operational context rather than a direct assumption that all conditions are identical to those at AIBD. In particular, the pressure thresholds and congestion classifications would require recalibration against local demand distributions, since the values used here are derived from AIBD’s own traffic composition. The relevance of this future direction lies in examining whether an integrated predictive operational intelligence approach can support arrival management beyond the original case while remaining sensitive to airport-specific processes and constraints.

In conclusion, the thesis demonstrates that AIBD’s arrival pressure can be addressed through a structured progression from operational explanation to prediction and decision support. The validated arrival dataset of 22,270 commercial records and the passenger- and baggage-pressure indicators derived from it provide the analytical foundation, identifying two recurring daily arrival banks at 04:00--06:00 and 19:00--21:00. Short-horizon forecasting anticipates passenger, baggage, and congestion conditions, with the deployed classifier achieving a ROC-AUC of 0.7914 once information leakage was diagnosed and corrected. AnyLogic simulation evaluates the operational consequences of these conditions across six scenarios, identifying baggage-reclaim capacity as the dominant constraint and demonstrating a 24 percent reduction in average passenger system time under integrated optimization. The Power BI Decision Support System translates the resulting evidence into manager-facing operational intelligence.

The central contribution is therefore not a single model, indicator, simulation scenario, or dashboard. It is the integration of these elements into an interpretable framework designed for a practical airport management problem. The framework supports preparation for arrival peaks, assessment of bottlenecks, comparison of operational scenarios, monitoring of interventions, and communication of decision-relevant information.

For AIBD, the approach aligns with the need to manage ongoing growth while maintaining operational efficiency and a passenger experience consistent with Teranga, reducing dependence on exclusively reactive management, and strengthening the capacity to anticipate pressure before it develops into severe congestion. The research consequently contributes at both academic and operational levels. Academically, it establishes an integrated, arrival-centric framework that connects proxy-based passenger- and baggage-pressure estimation and short-horizon forecasting with simulation and decision support in a West African airport context absent from the reviewed literature, while stating the limits that the aircraft-type proxy places on the demand layer itself. Operationally, it provides a structured basis for staffing preparation, baggage-capacity planning, queue-management strategies, intervention timing, and multi-stakeholder coordination. Within the limitations identified, the framework demonstrates how arrival data can be transformed into proactive operational intelligence and provides a foundation for the future extensions defined in this chapter.